Existuje mnoho způsobů, kterak v teorii množin vyrobit přirozená čísla. My se
přidržíme konstrukce, která přímo plyne z \hyperref[eq:axiom-nekonecna]{axiomu
nekonečna} a je naprosto názorná a dokonale neaxiomatická. Tedy, nestojí na
žádných výrocích, jež označíme za pravdivé, nýbrž je umožněna již poznanými
axiomy teorie množin.

V \hyperref[ssec:axiomy-vydeleni]{pododdíle o axiomech vydělení} jsme uvedli, že
jedním z jejich důsledků je existence prázdné množiny $0$. Stojí za uvědomění,
že $0$ je na první pohled vlastně \textbf{ta jediná} množina, o které můžeme
naprosto jednoznačně tvrdit, \textbf{že existuje}. Totiž,
\hyperref[eq:axiom-existence]{axiom existence} zaručuje existenci \emph{nějaké}
množiny $x$ a z \hyperref[eq:axiom-vydeleni]{axiomů vydělení} plyne, že existuje
i $x \setminus x$. Nic ale nezaručuje, že $x$ není sama množina $0$. Tedy, nic
až na \hyperref[eq:axiom-potence]{axiom potence} a již zmíněný
\hyperref[eq:axiom-nekonecna]{axiom nekonečna}. Díky těmto dvěma budeme moci
zkonstruovat přirozená čísla.

Nejprve ale poskytneme intuitivní představu této konstrukce: číslo $0$ je
označením pro prázdnou množinu právě pro to, že je množinou s \emph{nula}
prvky. Podobně, číslo $1$ bude označením pro množinu s \emph{jedním} prvkem,
číslo $2$ označením pro množinu se \emph{dvěma} prvky atd. Opakujeme,
\textbf{vše} v teorii množin \textbf{je množina}, tedy i čísla jsou množiny.

Kde ale ty další prvky vezmeme, když máme jen prázdnou množinu $0$? Přeci z
\hyperref[eq:axiom-potence]{axiomu potence}, který říká, že všechny podmnožiny
množiny $0$ tvoří množinu. Jistě, $0$ má jedinou podmnožinu, a to sebe sama, a
tedy $\mathcal{P}(0) = \{0\}$. Ale množina $\{0\}$ je množina s jedním prvkem!
Označíme ji jako $1$. Číslo $1$ je tedy množinou s jediným prvkem, kterým je
$0$.
